%=========================================================
\chapter{Requerimientos de usuario}
\label{cap:reqUsr}

    En este capitulo se presentan los requerimientos de usuario


%---------------------------------------------------------
\section{Requerimientos de usuario}

\begin{longtable}{|l|p{3cm}|p{7cm}|c|c|}
\caption{Requerimientos funcionales del sistema de gestión de hotel de perros.}
\label{tbl:reqFunc}\\
\hline
ID & Nombre & Descripción & Prioridad & Estado \\
\hline
\endfirsthead

\hline
ID & Nombre & Descripción & Prioridad & Estado \\
\hline
\endhead

\hline
\multicolumn{5}{r}{\footnotesize Continúa en la siguiente página} \\
\endfoot

\hline
\endlastfoot

		\FRitem{RU1}{Gestión de Empleados}{El usuario requiere llevar un registro actualizado de todos los empleados del hotel, incluyendo sus datos personales, información laboral (puesto, salario, fecha de contratación), datos de contacto y estado (activo/inactivo). El sistema debe permitir crear, modificar, consultar y dar de baja empleados.}{1}{\DONE}
		
		\FRitem{RU2}{Gestión de Propietarios}{El usuario requiere llevar un registro completo de los propietarios de mascotas, incluyendo datos personales, múltiples teléfonos de contacto, direcciones y su historial de reservaciones. El sistema debe facilitar la búsqueda rápida y actualización de información de contacto.}{1}{\DONE}
		
		\FRitem{RU3}{Gestión de Mascotas}{El usuario requiere mantener un registro detallado de cada mascota hospedada, incluyendo datos de identificación (nombre, especie, raza, color, sexo, fecha de nacimiento), historial médico completo (vacunaciones, desparasitaciones, enfermedades, alergias), documentos asociados y fotografías.}{1}{\DONE}
		
		\FRitem{RU4}{Sistema de Reservaciones}{El usuario requiere gestionar las reservaciones de hospedaje con un calendario visual interactivo que muestre la disponibilidad de habitaciones, permita crear nuevas reservas, asignar habitaciones, controlar check-in/check-out y llevar el seguimiento del estado de cada reservación.}{1}{\DONE}
		
		\FRitem{RU5}{Control de Habitaciones}{El usuario requiere llevar un registro actualizado de todas las habitaciones disponibles en el hotel, su capacidad, tipo, características especiales, estado de mantenimiento y disponibilidad para reservaciones.}{1}{\DONE}
		
		\FRitem{RU6}{Gestión de Servicios Adicionales}{El usuario requiere ofrecer y gestionar servicios complementarios al hospedaje tales como atención veterinaria, peluquería canina, entrenamiento y otros servicios especializados, con control de horarios, empleados asignados y tarifas.}{2}{\DONE}
		
		\FRitem{RU7}{Programación de Citas}{El usuario requiere un sistema de agendamiento para servicios especializados que permita programar citas, asignar empleados calificados, evitar conflictos de horario y llevar seguimiento del estado de cada cita (pendiente, realizada, cancelada).}{2}{\DONE}
		
		\FRitem{RU8}{Registro de Vacunaciones}{El usuario requiere mantener un historial completo y actualizado de todas las vacunas aplicadas a cada mascota, incluyendo tipo de vacuna, fecha de aplicación, fecha de próxima dosis y observaciones del veterinario.}{1}{\DONE}
		
		\FRitem{RU9}{Gestión de Documentos}{El usuario requiere almacenar, visualizar y gestionar documentos digitales importantes de cada mascota como certificados veterinarios, cartillas de vacunación, resultados de análisis y otros documentos oficiales en formatos PDF e imagen.}{2}{\DONE}
		
		\FRitem{RU10}{Control de Pagos}{El usuario requiere llevar un registro detallado de todos los pagos realizados por los clientes, métodos de pago utilizados, montos, fechas y asociarlos con las reservaciones o servicios correspondientes para control financiero.}{1}{\TODO}
		
		\FRitem{RU11}{Sistema de Autenticación}{El usuario requiere un sistema seguro de inicio de sesión que identifique a cada empleado y controle el acceso a las diferentes funcionalidades del sistema según su rol.}{1}{\DONE}
		
		\FRitem{RU12}{Control de Acceso por Roles}{El usuario requiere que el sistema restrinja las operaciones que cada empleado puede realizar según su puesto.}{1}{\DONE}
		
		\FRitem{RU13}{Historial de Hospedajes}{El usuario requiere consultar el historial completo de hospedajes de cada mascota para conocer fechas anteriores, habitaciones utilizadas, servicios contratados y cualquier incidente o nota relevante registrada.}{2}{\TODO}
		
		\FRitem{RU14}{Asignación de Empleados a Servicios}{El usuario requiere asignar empleados específicos a cada cita o servicio programado, considerando su especialización y disponibilidad, para garantizar la calidad del servicio.}{2}{\DONE}
		
		\FRitem{RU15}{Búsqueda y Filtros Avanzados}{El usuario requiere herramientas de búsqueda rápida y filtros para localizar fácilmente mascotas, propietarios, empleados y reservaciones usando criterios múltiples como nombre, fecha o estado.}{2}{\TODO}
		
		\FRitem{RU16}{Notificaciones de Vencimientos}{El usuario requiere recibir alertas automáticas sobre vacunas próximas a vencer, reservaciones próximas, pagos pendientes y otras fechas importantes.}{3}{\TODO}
		
		\FRitem{RU17}{Reportes de Ocupación}{El usuario requiere generar reportes periódicos sobre la ocupación del hotel, ingresos por período, servicios más solicitados y otros indicadores clave.}{3}{\TODO}
		
		\FRitem{RU18}{Gestión de Fotografías}{El usuario requiere almacenar y visualizar múltiples fotografías de cada mascota para su identificación visual y seguimiento durante el hospedaje.}{2}{\DONE}
		
		\FRitem{RU19}{Registro de Observaciones}{El usuario requiere documentar observaciones diarias sobre el comportamiento, alimentación, salud y cualquier incidente relevante de cada mascota durante su estancia.}{2}{\TODO}
		
		\FRitem{RU20}{Calendario Interactivo}{El usuario requiere una vista de calendario visual e interactiva que muestre todas las reservaciones, disponibilidad de habitaciones y citas de servicios.}{1}{\DONE}
		
		\FRitem{RU21}{Validación de Datos}{El usuario requiere que el sistema valide automáticamente la información ingresada para prevenir errores e inconsistencias.}{1}{\DONE}
		
		\FRitem{RU22}{Interfaz Responsiva}{El usuario requiere que el sistema sea accesible desde distintos dispositivos manteniendo funcionalidad y usabilidad.}{2}{\TODO}
		
		\FRitem{RU23}{Copia de Seguridad Automática}{El usuario requiere que el sistema realice respaldos automáticos periódicos de toda la información.}{1}{\TODO}
		
		\FRitem{RU24}{Auditoría de Cambios}{El usuario requiere llevar un registro de las modificaciones importantes realizadas en el sistema.}{3}{\TODO}
		
		\FRitem{RU25}{Exportación de Datos}{El usuario requiere exportar información del sistema en formatos estándar para análisis externo o respaldo.}{3}{\TODO}
        
        \FRitem{RU26}{Manejo de Habitaciones}{El usuario requiere administrar de forma detallada las habitaciones del hotel, incluyendo su alta, modificación y baja, así como la gestión de su capacidad, tipo, características especiales, estado de mantenimiento y disponibilidad para reservaciones.}{1}{\TODO}

        \FRitem{RU27}{Perfil de Empleado}{El usuario requiere que cada empleado cuente con un perfil individual que muestre su información personal, puesto, rol asignado, horarios, servicios que puede realizar y su historial de actividades dentro del sistema.}{2}{\TODO}

        \FRitem{RU28}{Manejo de Catálogos}{El usuario requiere gestionar los catálogos del sistema, tales como razas, tipos de servicios, tipos de habitaciones, estados, roles y otros datos parametrizables, permitiendo su creación, modificación y eliminación para mantener la flexibilidad del sistema.}{2}{\TODO}

		
	\end{longtable}

{\footnotesize\em Para leer correctamente esta tabla vea la leyenda en la Tabla~\ref{tbl:leyendaRF}.}



%=========================================================
\section{Requerimientos No Funcionales}
\label{sec:reqNoFunc}

En esta sección se presentan los requerimientos no funcionales del sistema, que especifican propiedades de calidad, restricciones y atributos que el sistema debe cumplir para garantizar su correcto funcionamiento, seguridad, rendimiento y mantenibilidad.

\begin{NFRequerimientos}
	\NFRitem{RNF1}{Seguridad}{
		El sistema debe garantizar que únicamente usuarios autenticados y autorizados puedan acceder a las funcionalidades del sistema mediante mecanismos robustos de autenticación y control de acceso basado en roles.
	}{
		Implementar autenticación mediante contraseñas encriptadas con bcrypt, control de acceso basado en roles (RBAC), sesiones con tokens JWT que expiren tras 7 días de inactividad, validación de permisos en cada operación sensible, y uso de HTTPS para comunicaciones.
	}{Crítica}
	
	\NFRitem{RNF2}{Rendimiento}{
		El sistema debe responder de forma ágil y eficiente ante las solicitudes del usuario, ejecutando operaciones en tiempos aceptables para no afectar la experiencia de usuario.
	}{
		Optimizar consultas a base de datos mediante índices apropiados, implementar caché para datos frecuentemente consultados, minimizar procesamiento innecesario en el cliente, establecer timeouts adecuados. Tiempo objetivo: consultas $<$ 1 segundo, operaciones de registro $<$ 2 segundos.
	}{Alta}
	
	\NFRitem{RNF3}{Escalabilidad}{
		El sistema debe permitir el crecimiento en volumen de datos y número de usuarios sin afectar su estabilidad ni requerir rediseños arquitectónicos mayores. Debe estar diseñado de forma modular y desacoplada.
	}{
		Utilizar arquitectura en capas claramente separadas (presentación, lógica de negocio, acceso a datos), aplicar principios SOLID para facilitar extensión, diseñar módulos independientes con interfaces bien definidas, implementar API RESTful stateless para permitir escalamiento horizontal.
	}{Alta}
	
	\NFRitem{RNF4}{Usabilidad}{
		El sistema debe proporcionar una interfaz clara, intuitiva y fácil de utilizar donde los usuarios puedan completar sus tareas de manera sencilla, con retroalimentación clara y prevención de errores.
	}{
		Diseñar interfaz consistente en todas las pantallas, implementar validación en tiempo real de formularios, proporcionar mensajes de error claros y específicos en español, usar retroalimentación visual para todas las acciones del usuario (loading states, confirmaciones), aplicar principios de diseño UX modernos.
	}{Alta}
	
	\NFRitem{RNF5}{Disponibilidad}{
		El sistema debe estar disponible y funcionar correctamente durante su horario operativo, manejando fallos de forma controlada y evitando la pérdida de información.
	}{
		Implementar manejo robusto de excepciones en todas las capas, registrar errores en logs para análisis posterior, realizar respaldos automáticos periódicos de datos críticos, diseñar mecanismos de recuperación ante fallos, establecer monitoreo del estado del sistema.
	}{Alta}
	
	\NFRitem{RNF6}{Mantenibilidad}{
		El sistema debe facilitar la comprensión, modificación y evolución del código por parte del equipo de desarrollo mediante código claro, estructurado y siguiendo buenas prácticas de programación.
	}{
		Aplicar convenciones de nomenclatura consistentes (camelCase para variables/métodos, PascalCase para clases/componentes), documentar código y funciones complejas mediante comentarios descriptivos, estructurar el proyecto de forma lógica y modular, seguir estándares de codificación establecidos (ESLint, Prettier), realizar revisiones de código en equipo.
	}{Media}
	
	\NFRitem{RNF7}{Compatibilidad}{
		El sistema debe funcionar correctamente en distintos entornos de ejecución, navegadores web modernos y sistemas operativos, sin requerir configuraciones especiales del usuario final.
	}{
		Utilizar tecnologías web estándar (HTML5, CSS3, JavaScript ES6+), probar funcionamiento en navegadores principales (Chrome, Firefox, Safari, Edge versiones recientes), diseñar para diferentes resoluciones de pantalla (responsive design), facilitar despliegue en diferentes entornos (desarrollo, staging, producción) mediante variables de configuración.
	}{Media}
	
	\NFRitem{RNF8}{Integridad de Datos}{
		El sistema debe garantizar que la información almacenada sea correcta, consistente y coherente, validando los datos antes de ser almacenados y manteniéndolos íntegros durante todo su ciclo de vida.
	}{
		Implementar validación de datos en múltiples capas (cliente y servidor), establecer restricciones de integridad en base de datos (claves foráneas, valores únicos, NOT NULL, CHECK constraints), validar tipos y formatos de datos, implementar transacciones ACID para operaciones críticas, proporcionar mensajes claros de error de validación al usuario.
	}{Crítica}
	
	\NFRitem{RNF9}{Respaldo y Recuperación}{
		El sistema debe contar con mecanismos automatizados de respaldo de información y procedimientos documentados para recuperación ante desastres o pérdida de datos.
	}{
		Implementar respaldos automáticos diarios de la base de datos completa, mantener respaldos incrementales cada 6 horas, almacenar respaldos en ubicación separada del servidor principal, documentar procedimientos de restauración, probar periódicamente la recuperación de respaldos.
	}{Alta}
	
	\NFRitem{RNF10}{Auditabilidad}{
		El sistema debe mantener un registro de las operaciones críticas realizadas, identificando quién realizó cada acción, cuándo y sobre qué datos, para fines de auditoría y rastreo de problemas.
	}{
		Registrar en logs todas las operaciones de creación, modificación y eliminación de datos sensibles, incluir timestamp, usuario responsable y datos modificados en cada registro, implementar sistema de logs centralizado con niveles de severidad (ERROR, WARN, INFO), permitir consulta y filtrado de logs por administradores.
	}{Media}
	
	\NFRitem{RNF11}{Privacidad}{
		El sistema debe proteger la información personal de los usuarios y cumplir con regulaciones de protección de datos aplicables, limitando el acceso a información sensible solo a usuarios autorizados.
	}{
		Encriptar contraseñas usando algoritmos fuertes (bcrypt con salt), no almacenar información de tarjetas de crédito en el sistema, implementar políticas de privacidad claras, permitir a propietarios consultar y modificar únicamente sus propios datos, limitar acceso de empleados a información estrictamente necesaria para su rol.
	}{Crítica}
	
	\NFRitem{RNF12}{Confiabilidad}{
		El sistema debe operar de manera predecible y consistente, minimizando errores, comportamientos inesperados y caídas del sistema durante su operación normal.
	}{
		Implementar pruebas unitarias para componentes críticos, realizar pruebas de integración entre módulos, validar todas las entradas de usuario antes de procesarlas, manejar casos excepcionales de forma explícita (conexiones de red fallidas, datos faltantes), implementar circuit breakers para servicios externos.
	}{Alta}
	
	\NFRitem{RNF13}{Portabilidad}{
		El sistema debe poder ser desplegado en diferentes entornos de hosting y proveedores cloud sin requerir modificaciones significativas en el código fuente.
	}{
		Utilizar contenedores Docker para empaquetar la aplicación con todas sus dependencias, externalizar configuración mediante variables de ambiente (.env), documentar requerimientos de infraestructura (versiones de Node.js, PostgreSQL), evitar dependencias de servicios propietarios de proveedores específicos.
	}{Media}
	
	\NFRitem{RNF14}{Documentación}{
		El sistema debe estar acompañado de documentación técnica completa y actualizada que facilite su instalación, configuración, uso y mantenimiento por parte de desarrolladores y administradores.
	}{
		Mantener archivo README.md con instrucciones de instalación y ejecución, documentar API REST con descripción de endpoints, parámetros y respuestas, incluir ejemplos de uso de funcionalidades principales, documentar modelo de base de datos con diagrama ER y descripción de tablas, mantener changelog con registro de cambios por versión.
	}{Media}
	
	\NFRitem{RNF15}{Internacionalización}{
		El sistema debe estar preparado para soportar múltiples idiomas en el futuro, separando textos de la interfaz de usuario del código de la aplicación.
	}{
		Externalizar todos los textos mostrados al usuario en archivos de configuración separados (archivos de traducción), utilizar biblioteca de i18n (internacionalización) compatible con React, codificar cadenas de texto en UTF-8 para soportar caracteres especiales, diseñar interfaz considerando expansión/contracción de textos en diferentes idiomas.
	}{Baja}
\end{NFRequerimientos}

{\footnotesize\em Los requerimientos no funcionales complementan los funcionales estableciendo propiedades de calidad que el sistema debe cumplir.}
